% GENERATED FILE. Edit canonical lessons, then run npm run generate:books.
% canonical-source-hash: fnv1a64:2741f0abc2290535
% canonical-lessons: AR-C323-nashaa, AR-C323-taraara, AR-C323-masqat-ras, AR-C323-siba, AR-C323-kuhula

\chapter{Growing up and growing older}
\label{ch:ar-growing-up-and-growing-older}
\hlchaptermodality{323}

\begin{chapteropening}
\textbf{By the end of this chapter:} \emph{I can say to be reared, to be brought up, a birthplace, boyhood and middle age.}

Complete the last lesson of chapter 323: I can say to be reared, to be brought up, a birthplace, boyhood and middle age.
\end{chapteropening}

\section[\texorpdfstring{\ar{نشأ} (nashaʾa)}{نشأ (nashaʾa)}]{\ar{نشأ} (nashaʾa) — to be reared, to grow up; to arise}

\label{lesson:AR-C323-nashaa}



\begin{quote}
Before the new one: say the Arabic for to besiege, then the Arabic for to defend.
\end{quote}

\subsection*{You'll want to know: \ar{نشأ}}
\textbf{\ar{نشأ}} — \emph{nashaʾa} — \textquotedblleft{}to be reared, to grow up; to arise\textquotedblright{}.

\begin{quote}
\ar{نشأت} \ar{أمي} \ar{في} \ar{قرية} \ar{صغيرة}.
\end{quote}

\begin{quote}
\emph{nashaʾat ummī fī qarya ṣaghīra.}
\end{quote}

\begin{quote}
My mother grew up in a small village.
\end{quote}

Present: \ar{ينشأ}, \emph{yanshaʾ}.

Of a problem or an idea: to arise.

\begin{cousinweb}
\textbf{\ar{نشأ}} (\emph{nashaʾa}) is built on the native Arabic root \emph{n-sh-ʾ} (Lane), \textquotedblleft{}to come into being; to grow up\textquotedblright{}; \emph{nashʾa} is upbringing.
\end{cousinweb}

\begin{grammarlens}[title={What you've built}]
One more word for telling what happened.
\end{grammarlens}

\subsection*{Your turn}
\begin{itemize}
\raggedright
  \item \emph{Say it:} \emph{nashaʾa}
  \item \emph{Say it:} \emph{nashaʾa}, once more
  \item \emph{Say it:} \emph{nashaʾat ummī fī qarya ṣaghīra.} — My mother grew up in a small village.
  \item \emph{Recall:} say the Arabic for to besiege, then the Arabic for to defend, then say \emph{nashaʾa} again
\end{itemize}

\begin{tcolorbox}[breakable,colback=teal!4,colframe=teal!35!black,title={Before you move on}]
What does \ar{نشأ} mean? (\textquotedblleft{}To be reared, to grow up; to arise\textquotedblright{}.) Say it once more.
\end{tcolorbox}
\section[\texorpdfstring{\ar{ترعرع} (taraʿraʿa)}{ترعرع (taraʿraʿa)}]{\ar{ترعرع} (taraʿraʿa) — to be brought up}

\label{lesson:AR-C323-taraara}



\begin{quote}
Before the new one: say the Arabic for to defend, then the Arabic for to be reared.
\end{quote}

\subsection*{You'll want to know: \ar{ترعرع}}
\textbf{\ar{ترعرع}} — \emph{taraʿraʿa} — \textquotedblleft{}to be brought up\textquotedblright{}.

\begin{quote}
\ar{ترعرع} \ar{أبي} \ar{قرب} \ar{البحر}.
\end{quote}

\begin{quote}
\emph{taraʿraʿa abī qurba al-baḥr.}
\end{quote}

\begin{quote}
My father was brought up near the sea.
\end{quote}

Present: \ar{يترعرع}, \emph{yataraʿraʿ}.

A literary word beside \ar{نشأ}.

\begin{cousinweb}
\textbf{\ar{ترعرع}} (\emph{taraʿraʿa}) is built on the native Arabic root \emph{r-ʿ-ʿ} (Lane), whose doubled verb \emph{raʿraʿa} is \textquotedblleft{}to make a child grow\textquotedblright{}. \emph{Taraʿraʿa}, its reflexive, is to grow up.
\end{cousinweb}

\begin{grammarlens}[title={What you've built}]
One more word for telling what happened.
\end{grammarlens}

\subsection*{Your turn}
\begin{itemize}
\raggedright
  \item \emph{Say it:} \emph{taraʿraʿa}
  \item \emph{Say it:} \emph{taraʿraʿa}, once more
  \item \emph{Say it:} \emph{taraʿraʿa abī qurba al-baḥr.} — My father was brought up near the sea.
  \item \emph{Recall:} say the Arabic for to defend, then the Arabic for to be reared, then say \emph{taraʿraʿa} again
\end{itemize}

\begin{tcolorbox}[breakable,colback=teal!4,colframe=teal!35!black,title={Before you move on}]
What does \ar{ترعرع} mean? (\textquotedblleft{}To be brought up\textquotedblright{}.) Say it once more.
\end{tcolorbox}
\section[\texorpdfstring{\ar{مسقط} \ar{رأس} (masqaṭ raʾs)}{مسقط رأس (masqaṭ raʾs)}]{\ar{مسقط} \ar{رأس} (masqaṭ raʾs) — a birthplace, a home town}

\label{lesson:AR-C323-masqat-ras}



\begin{quote}
Before the new one: say the Arabic for to be reared, then the Arabic for to be brought up.
\end{quote}

\subsection*{You'll want to know: \ar{مسقط} \ar{رأس}}
\textbf{\ar{مسقط} \ar{رأس}} — \emph{masqaṭ raʾs} — \textquotedblleft{}a birthplace, a home town\textquotedblright{}.

\begin{quote}
\ar{زرت} \ar{مسقط} \ar{رأس} \ar{أبي} \ar{في} \ar{الصيف}.
\end{quote}

\begin{quote}
\emph{zurtu masqaṭ raʾs abī fī aṣ-ṣayf.}
\end{quote}

\begin{quote}
I visited my father’s home town in the summer.
\end{quote}

Literally the place where the head fell: where one was born.

\begin{cousinweb}
\textbf{\ar{مسقط} \ar{رأس}} (\emph{masqaṭ raʾs}) is \textquotedblleft{}the place where the head fell\textquotedblright{}: \emph{masqaṭ} is built on the native Arabic root \emph{s-q-ṭ}, \textquotedblleft{}to fall\textquotedblright{} (Lane), and \emph{raʾs}, a head, is an inherited Semitic word, beside Hebrew \emph{rōsh} and Aramaic \emph{rēshā} (EtymDB).
\end{cousinweb}

\begin{grammarlens}[title={What you've built}]
One more word for telling what happened.
\end{grammarlens}

\subsection*{Your turn}
\begin{itemize}
\raggedright
  \item \emph{Say it:} \emph{masqaṭ raʾs}
  \item \emph{Say it:} \emph{masqaṭ raʾs}, once more
  \item \emph{Say it:} \emph{zurtu masqaṭ raʾs abī fī aṣ-ṣayf.} — I visited my father’s home town in the summer.
  \item \emph{Recall:} say the Arabic for to be reared, then the Arabic for to be brought up, then say \emph{masqaṭ raʾs} again
\end{itemize}

\begin{tcolorbox}[breakable,colback=teal!4,colframe=teal!35!black,title={Before you move on}]
What does \ar{مسقط} \ar{رأس} mean? (\textquotedblleft{}A birthplace, a home town\textquotedblright{}.) Say it once more.
\end{tcolorbox}
\section[\texorpdfstring{\ar{صبا} (ṣibā)}{صبا (ṣibā)}]{\ar{صبا} (ṣibā) — boyhood, early youth}

\label{lesson:AR-C323-siba}



\begin{quote}
Before the new one: say the Arabic for to be brought up, then the Arabic for a birthplace.
\end{quote}

\subsection*{You'll want to know: \ar{صبا}}
\textbf{\ar{صبا}} — \emph{ṣibā} — \textquotedblleft{}boyhood, early youth\textquotedblright{}.

\begin{quote}
\ar{قضى} \ar{جدي} \ar{صباه} \ar{في} \ar{الريف}.
\end{quote}

\begin{quote}
\emph{qaḍā jaddī ṣibāhu fī ar-rīf.}
\end{quote}

\begin{quote}
My grandfather spent his boyhood in the countryside.
\end{quote}

A literary word; the east wind is also \ar{صبا}, \emph{ṣabā}.

\begin{cousinweb}
\textbf{\ar{صبا}} (\emph{ṣibā}) is built on the native Arabic root \emph{ṣ-b-w} (Lane), \textquotedblleft{}to incline, to be youthful\textquotedblright{}, the root of \emph{ṣabiyy}, a boy. \emph{Ṣibā} is the time of being one.
\end{cousinweb}

\begin{grammarlens}[title={What you've built}]
One more word for telling what happened.
\end{grammarlens}

\subsection*{Your turn}
\begin{itemize}
\raggedright
  \item \emph{Say it:} \emph{ṣibā}
  \item \emph{Say it:} \emph{ṣibā}, once more
  \item \emph{Say it:} \emph{qaḍā jaddī ṣibāhu fī ar-rīf.} — My grandfather spent his boyhood in the countryside.
  \item \emph{Recall:} say the Arabic for to be brought up, then the Arabic for a birthplace, then say \emph{ṣibā} again
\end{itemize}

\begin{tcolorbox}[breakable,colback=teal!4,colframe=teal!35!black,title={Before you move on}]
What does \ar{صبا} mean? (\textquotedblleft{}Boyhood, early youth\textquotedblright{}.) Say it once more.
\end{tcolorbox}
\section[\texorpdfstring{\ar{كهولة} (kuhūla)}{كهولة (kuhūla)}]{\ar{كهولة} (kuhūla) — middle age}

\label{lesson:AR-C323-kuhula}



\begin{quote}
Before the new one: say the Arabic for a birthplace, then the Arabic for boyhood.
\end{quote}

\subsection*{You'll want to know: \ar{كهولة}}
\textbf{\ar{كهولة}} — \emph{kuhūla} — \textquotedblleft{}middle age\textquotedblright{}.

\begin{quote}
\ar{بدأ} \ar{أبي} \ar{الرسم} \ar{في} \ar{كهولته}.
\end{quote}

\begin{quote}
\emph{badaʾa abī ar-rasm fī kuhūlatihi.}
\end{quote}

\begin{quote}
My father took up painting in middle age.
\end{quote}

A man in middle age is \ar{كهل}.

\begin{cousinweb}
\textbf{\ar{كهولة}} (\emph{kuhūla}) is built on the native Arabic root \emph{k-h-l} (Lane), where \emph{kahl} is a man past youth whose hair is turning grey. \emph{Kuhūla} is that time of life.
\end{cousinweb}

\begin{grammarlens}[title={What you've built}]
One more word for telling what happened.
\end{grammarlens}

\subsection*{Your turn}
\begin{itemize}
\raggedright
  \item \emph{Say it:} \emph{kuhūla}
  \item \emph{Say it:} \emph{kuhūla}, once more
  \item \emph{Say it:} \emph{badaʾa abī ar-rasm fī kuhūlatihi.} — My father took up painting in middle age.
  \item \emph{Recall:} say the Arabic for a birthplace, then the Arabic for boyhood, then say \emph{kuhūla} again
\end{itemize}

\begin{tcolorbox}[breakable,colback=teal!4,colframe=teal!35!black,title={Before you move on}]
What does \ar{كهولة} mean? (\textquotedblleft{}Middle age\textquotedblright{}.) Say it once more.
\end{tcolorbox}
